% Shared LaTeX macros for this paper's reproduction package.
% Both papers \input this file.

% ----- model class shorthands -----
\newcommand{\Mone}{M\textsubscript{1}\xspace}      % flat
\newcommand{\Mtwo}{M\textsubscript{2}\xspace}      % hierarchical regime-aware
\newcommand{\LEAR}{LEAR\xspace}
\newcommand{\XGB}{XGBoost\xspace}

% ----- units -----
\newcommand{\jpykwh}{\unit{JPY/kWh}}
\newcommand{\maeunit}{\unit{JPY/kWh}}

% ----- numbers we cite repeatedly -----
% Primary headline numbers refer to the daily-rolling walk-forward
% (step = 3 days, n = 140 folds, full weekday coverage via gcd(3, 7) = 1).
% The earlier weekly-anchor design (step = 7 d, n = 60, Sunday-only test
% days) is retained internally as a comparator -- see Section 4 and
% Appendix C for the weekly-vs-daily robustness table.
\newcommand{\SOTA}{\num{1.816 \pm 0.92}}        % Kyushu M1/XGB+weather, daily n=140
\newcommand{\SOTAweekly}{\num{1.573 \pm 0.75}}  % Kyushu M1/XGB+weather, weekly n=60 (legacy)
\newcommand{\Spearman}{\rho = -0.733}           % cross-area FLOOR vs Δ MAE (daily)
\newcommand{\SpearmanP}{p = 0.025}
\newcommand{\SpearmanPre}{\rho = -0.767}        % FLOOR computed pre-evaluation 2022-01..2024-09-14
\newcommand{\SpearmanPreP}{p = 0.016}
\newcommand{\SpearmanTrain}{\rho = -0.683}      % FLOOR computed FY2024-only (training-window approx)
\newcommand{\SpearmanTrainP}{p = 0.042}
\newcommand{\CurtailFloorSpearman}{\rho = +0.833}  % cross-area curtailment-hours vs FLOOR-hours (8-area incl Shikoku+Kansai, 2026-06-17)
\newcommand{\CurtailFloorSpearmanP}{p = 0.010}
\newcommand{\CurtailFloorN}{n = 8}

% ----- Round-27 / Phase 26: operational-forecast reframe -----
% Within-day triple under a bid-close-admissible *operational-forecast*
% moderator (Phase 26.0b isolation on the committed full-sample DV): the
% mechanism survives admissibility; the lag-48 collapse to NS was
% measurement-error attenuation from a noisy moderator (midday corr 0.57 vs
% the histfc forecast's 0.89), not a look-ahead artifact.
\newcommand{\WithinDayBeta}{-0.236}            % within-day triple, histfc-forecast moderator
\newcommand{\WithinDayP}{p = 0.049}            % wild-cluster bootstrap, G=9 (full sample)
\newcommand{\WithinDayBetaContemp}{-0.342}     % contemporaneous-moderator comparison (reproduces committed -0.33)
\newcommand{\WithinDayFEbeta}{-0.012}          % within-day triple, histfc moderator, saturated area+date+hour FE (phase28.1 spec1)
\newcommand{\WithinDayFEp}{0.96}               % two-way (area x date) clustered asymptotic p under saturated FE
% Round-33 B3: reviewer-requested LITERAL slot-level PV-surplus intensity panel
% (phase30.1: |eps_w|-|eps_p| ~ realised PV-to-demand x midday + FE). Reported as a
% transparent null (FE-fragile), consistent with H1' and the moderated-triple collapse.
\newcommand{\IntensityNoFEbeta}{-0.228}        % no-FE slope (wild-cluster p<0.001), suggestive within-area intensity
\newcommand{\IntensityNoFEp}{<0.001}           % wild-cluster bootstrap two-sided p, no-FE
\newcommand{\IntensityFEbeta}{-0.157}          % saturated area+date+hour FE slope (attenuated)
\newcommand{\IntensityFEp}{0.30}               % two-way clustered p under saturated FE (n.s.)
% Round-30 Gate 4 (D-lite): observed area-vs-system price wedge during FLOOR episodes
% (phase29), grounding the asymmetric cost-ratio range in the high-FLOOR areas.
\newcommand{\FloorWedgeShikoku}{6.33}          % Shikoku floor-episode median (system - area) JPY/kWh
\newcommand{\FloorWedgeKyushu}{3.1}            % Kyushu floor-episode median (system - area) JPY/kWh
% Operational forecast (JMA-MSM histfc proxy) vs ERA5-lag48 proxy, clean
% same-sample head-to-head (Phase 26.1, 8 OCCTO areas): the cross-area
% concentration is a property of the WEAK proxy and washes out under a
% genuine operational forecast, which is ~7x more informative and broad.
\newcommand{\OpsCrossArea}{\rho = -0.29}       % operational cross-area gradient (improvement vs PV/peak), n.s.
\newcommand{\OpsCrossAreaP}{p = 0.49}
\newcommand{\ProxyCrossAreaClean}{\rho = +0.76}  % ERA5 proxy, clean sample, on IMPROVEMENT (=-Delta MAE), p=0.028; equiv rho=-0.76 on signed Delta MAE (sign convention: sec:per_area)
% Decision-loss economic value (Phase 26.2): the operational forecast has
% real floor-event decision value, unlike the battery-arbitrage null.
\newcommand{\OpsFloorFone}{0.28}               % operational floor-event F1 (price-only 0.13)
\newcommand{\PriceOnlyFloorFone}{0.13}
\newcommand{\OpsFloorAUC}{0.96}                % operational floor-event ROC-AUC (price-only 0.93)

% ----- Round-28 / Phase 27: threshold-crossing VoI + deepened decision value -----
% Floor-event decision value, deepened (Phase 26.2 Round-28): the operational
% forecast captures nearly all of the realised-weather benchmark's mean weather
% value (not a strict per-metric upper bound); the lagged proxy captures almost
% none. Base rate explains the high-AUC / modest-F1 gap.
\newcommand{\FloorBaseRate}{3.3\%}                 % fraction of slots at FLOOR (explains F1 vs AUC)
\newcommand{\OpsFloorPRAUC}{0.59}                  % operational floor-event PR-AUC (price-only 0.37)
\newcommand{\PriceOnlyFloorPRAUC}{0.37}
\newcommand{\OracleFloorPRAUC}{0.57}              % realised-weather oracle PR-AUC (operational ~= oracle)
\newcommand{\OracleFloorFone}{0.29}              % realised-weather oracle floor F1 (operational 0.28)
\newcommand{\OpsFloorBrier}{0.020}               % operational Brier (price-only 0.025), cross-fitted isotonic
\newcommand{\PriceOnlyFloorBrier}{0.025}
\newcommand{\OpsFloorFoneCF}{0.54}               % operational floor F1 at cross-fitted decision threshold
\newcommand{\PriceOnlyFloorFoneCF}{0.43}
\newcommand{\OpsCurtailCoincidence}{83\%}        % operational-predicted floors coincident with curtailment
\newcommand{\PriceOnlyCurtailCoincidence}{53\%}
% 4-information-set decomposition of mean weather value (yen/kWh; Phase 27.1):
% proxy -> operational -> realised-weather oracle. Operational ~= oracle.
\newcommand{\ImpProxy}{0.034}                    % ERA5 lag-48 proxy mean MAE improvement
\newcommand{\ImpOps}{0.239}                      % operational JMA-MSM mean MAE improvement
\newcommand{\ImpOracle}{0.268}                   % realised-weather oracle mean MAE improvement
\newcommand{\OpsCaptureOracle}{89\%}             % operational / oracle mean weather value
\newcommand{\ProxyCaptureOracle}{13\%}           % proxy / oracle mean weather value
% Forecast skill: SkillProxy/SkillHistfc are the midday corr(forecast, realised shortwave) = sqrt(tau); model defines tau = corr^2 (Phase 27.0).
\newcommand{\SkillProxy}{0.57}                   % ERA5 lag-48 proxy skill: midday forecast-vs-realised corr (= sqrt(tau); model tau=corr^2)
\newcommand{\SkillHistfc}{0.89}                  % operational JMA-MSM skill: midday forecast-vs-realised corr (= sqrt(tau))
% Round-37 Phase 3 (review #17): CHANNEL-resolved forecast skill = corr(forecast,
% realised ERA5), cross-area mean, MIDDAY (09-15 JST). proxy=ERA5 lag-48, gpv=exact-
% issuance JMA-GPV. Grounds H1' channel breadth (eq:voi_decomp): the proxy is
% narrowband (skilled on diurnally-persistent temperature+shortwave, poor on cloud+
% wind); the operational forecast's skill GAIN is largest on the kappa-free broad
% channels (wind+cloud), which is why it washes out the PV-floor concentration.
% Source: phase42_channel_skill.json.
\newcommand{\ChanProxySW}{0.50}                  % proxy shortwave skill (midday corr)
\newcommand{\ChanGpvSW}{0.83}                    % GPV shortwave skill
\newcommand{\ChanProxyTemp}{0.95}                % proxy temperature skill (near-saturated, diurnal persistence)
\newcommand{\ChanGpvTemp}{0.99}                  % GPV temperature skill
\newcommand{\ChanProxyCloud}{0.17}               % proxy cloud skill (narrowband: poor)
\newcommand{\ChanGpvCloud}{0.60}                 % GPV cloud skill
\newcommand{\ChanProxyWind}{0.29}                % proxy wind skill (narrowband: poor)
\newcommand{\ChanGpvWind}{0.73}                  % GPV wind skill
\newcommand{\ChanGainWind}{0.44}                 % GPV skill GAIN over proxy on wind (= displayed 0.73-0.29; JSON 0.446); kappa-free channel
\newcommand{\ChanGainCloud}{0.43}                % GPV skill GAIN over proxy on cloud (= displayed 0.60-0.17; JSON 0.433); PV/floor (solar) channel
% ----- Round-31 / Phase 31: exact-issuance GPV head-to-head + real-yen decision value -----
% Full-period TRUE exact-issuance JMA-GPV forecast (Kyoto RISH, D-1 00UTC MSM run,
% all 426 days) re-run head-to-head on the daylight identification window (FH16-33).
% It reproduces the Historical-Forecast RECONSTRUCTION: large broad value, washed-out
% PV-floor concentration -> the Table 8 washout is genuine forecast skill, not a
% reconstruction artefact (review #1). gpv 0.314 ~= histfc 0.297 on this daytime sample.
\newcommand{\GpvExactImp}{0.314}                 % exact-issuance GPV mean MAE improvement (daytime same-sample)
\newcommand{\HistfcDayImp}{0.297}                % reconstruction mean improvement, identical daytime sample
\newcommand{\GpvExactConc}{\rho = -0.26}         % exact-issuance GPV cross-area gradient vs PV/peak (washout, n.s.)
\newcommand{\GpvExactConcP}{p = 0.53}
% Round-33 Phase A: ALL-HOURS (00-23 JST) exact-issuance head-to-head, the headline
% (phase35_allhours_headtohead.json, n=44,928, 8 areas). gpv/proxy ~7.9x, washout holds;
% proxy concentration now SIGNIFICANT all-hours (+0.83, p=0.01), removing the daylight caveat.
\newcommand{\ProxyAllHoursImp}{0.032}            % ERA5 proxy mean MAE improvement (all-hours same-sample)
\newcommand{\HistfcAllHoursImp}{0.240}           % operational reconstruction mean improvement (all-hours)
\newcommand{\GpvAllHoursImp}{0.253}              % exact-issuance GPV mean MAE improvement (all-hours)
\newcommand{\ProxyAllHoursConc}{+0.83}           % proxy cross-area gradient vs PV/peak (imp convention; significant)
\newcommand{\ProxyAllHoursConcP}{0.01}
\newcommand{\HistfcAllHoursConc}{-0.29}          % operational reconstruction gradient (washout, n.s.)
\newcommand{\HistfcAllHoursConcP}{0.49}
\newcommand{\GpvAllHoursConc}{\rho = -0.14}      % exact-issuance GPV gradient vs PV/peak (washout, n.s.)
\newcommand{\GpvAllHoursConcP}{p = 0.74}
% Round-37 Phase 2 (review #17): 9-area point-forecast head-to-head (adds Chubu;
% feasible because the point-forecast comparison needs only price+weather, not OCCTO
% supply data -- decision-value suite stays 8-area). Confirms the H1' pattern is not
% an 8-area artefact. Source: phase41_allhours_headtohead_9area.json (n=50,544).
\newcommand{\ProxyAllHoursConcNine}{+0.72}       % 9-area proxy gradient vs PV/peak (imp convention; still significant)
\newcommand{\ProxyAllHoursConcNineP}{0.03}
\newcommand{\HistfcAllHoursConcNine}{-0.17}      % 9-area reconstruction gradient (washout, n.s.)
\newcommand{\HistfcAllHoursConcNineP}{0.67}
\newcommand{\GpvAllHoursConcNine}{-0.02}         % 9-area exact-issuance GPV gradient (washout, n.s.)
\newcommand{\GpvAllHoursConcNineP}{0.97}
% Round-34 Phase 3: controlled skill-degradation sweep (phase36_skill_degradation.json).
% A SINGLE-source comparative static -- the all-hours exact-issuance JMA-GPV forecast
% degraded by increasing Gaussian noise -- traces mean weather value monotonically up
% with forecast skill (corr with realised). The concentration half is a signal-STRUCTURE
% effect (channel-recovery ratio), not a scalar-skill one: isotropic noise never develops
% the proxy's significant PV-floor concentration (|rho|<=0.26, n.s. at every level).
\newcommand{\SkillCurveHiSkill}{0.841}           % forecast skill (corr) at the pristine-forecast level
\newcommand{\SkillCurveHiVal}{0.258}             % mean weather value (JPY/kWh) at the pristine-forecast level
\newcommand{\SkillCurveLoSkill}{0.175}           % forecast skill (corr) at the most-degraded level
\newcommand{\SkillCurveLoVal}{0.026}             % mean weather value (JPY/kWh) at the most-degraded level
% Round-35 Phase B: area-month panel (phase37_areamonth_panel.json) escaping the n=9
% cross-section. DV = mean fold weather Delta-MAE per (area, month); treatment = within-area
% time-varying FLOOR rate. Negative beta = higher-FLOOR months carry more weather value.
% Significant under area FE (within-area temporal identification, wild-cluster G=9),
% attenuating under fully-saturated area+month FE (honest, same pattern as the slot panel).
\newcommand{\PanelMonthN}{135}                   % (area, month) cells (9 areas x 15 months, >=3 folds)
\newcommand{\PanelMonthBetaPooled}{-1.786}       % spec0 pooled beta on FLOOR rate
\newcommand{\PanelMonthBeta}{-1.917}             % spec1 area-FE beta on FLOOR (within-area temporal, headline)
\newcommand{\PanelMonthWildP}{0.004}             % spec1 wild-cluster bootstrap p (G=9)
\newcommand{\PanelMonthBetaSatFE}{-1.330}        % spec2 area+month-FE beta (attenuated)
\newcommand{\PanelMonthSatP}{0.308}              % spec2 wild-cluster p (n.s. under full saturation)
% Round-35 Phase E: external-validity exploratory on EPEX-DE/LU (phase38_epex_de_exploratory.json).
% Public keyless data (SMARD price/PV/load + Open-Meteo). Weather value replicates in DE, but the
% floor-concentration analogue (vs negative-price exposure) is only weakly positive and n.s. -- DE's
% oversupply boundary is softer and wind-co-driven, not a sharp PV floor (gate: exploratory, not headline).
\newcommand{\DEWeatherValue}{0.757}              % DE mean weather improvement, EUR/MWh (2024)
\newcommand{\DENdays}{239}                       % DE walk-forward evaluation days
\newcommand{\DEConcRho}{0.098}                   % Spearman(daily weather improvement, negative-price share)
\newcommand{\DEConcP}{0.133}                     % p-value (n.s.) -- concentration weak at the softer boundary
% Real-yen floor-event decision value (Phase 31c): per-slot floor decisions joined to
% REAL JEPX area imbalance-settlement (shortage) prices from the 9 TD operators'
% disclosure. Replaces the normalized-illustrative ~3.0M scaling (review #3).
\newcommand{\RealFloorWedge}{6.34} % RETIRED (phase57): one-sided metric, wrong per-MW scaling; do NOT cite               % median real floor over-cost (shortage imbalance - spot) JPY/kWh
\newcommand{\RealDecisionCostPriceOnly}{1.89} % RETIRED (phase57): one-sided metric, wrong per-MW scaling; do NOT cite    % price-only real imbalance exposure, M JPY/MW/yr
\newcommand{\RealDecisionCostOps}{1.71} % RETIRED (phase57): one-sided metric, wrong per-MW scaling; do NOT cite          % operational real imbalance exposure, M JPY/MW/yr
\newcommand{\OpsRealValueK}{177} % RETIRED (phase57): one-sided metric, wrong per-MW scaling; do NOT cite                 % operational forecast real decision value, thousand JPY/MW/yr
% Honest enriched-baseline robustness (Phase 32, P5): enriched PUBLIC baseline
% (price + system + neighbouring-area prices + OCCTO supply-mix). The lag-48 supply
% variant is an OVB-robustness baseline (supply actuals at lag-48 are NOT strictly
% bid-close-admissible for slots>10:00); weather value is NOT a weak-baseline
% artefact (review #7/#2).
\newcommand{\EnrichBaselineGain}{0.057}          % mean MAE gain of lag-48-supply enriched baseline over price-only (8/8) -> strong
\newcommand{\OpsOverEnriched}{0.19}              % operational-forecast mean MAE improvement OVER the lag-48-supply enriched baseline (DM-sig 8/8)
\newcommand{\OpsOverEnrichedNsig}{8}             % areas where operational fc beats enriched baseline at DM p<0.05 (of 8)
\newcommand{\OpsOverEnrichedConc}{+0.05}         % cross-area Spearman of operational improvement vs PV/peak (broad, n.s.)
\newcommand{\ProxyOverEnriched}{0.015}           % ERA5 proxy mean MAE improvement over the enriched baseline (information-poor)
% Round-37 Phase 1 (review #17 must-fix): GENUINELY bid-close-admissible enriched
% baseline -- supply mix moved to lag-96 (D-2 same-clock, observed by D-1 10:00 for
% every slot), prices stay at lag-48 (day-ahead prices for D-1 publish at the D-2
% auction). Operational forecast survives this strictly-admissible baseline even
% more strongly than the lag-48 variant. Source: phase40_admissible_supply_lag96.json.
\newcommand{\EnrichBaselineGainAdm}{0.032}       % mean MAE gain of admissible (lag-96-supply) enriched baseline over price-only (8/8) -> strong
\newcommand{\OpsOverEnrichedAdm}{0.216}          % operational-forecast mean MAE improvement OVER the admissible enriched baseline (DM-sig 8/8)
\newcommand{\OpsOverEnrichedAdmNsig}{8}          % areas where operational fc beats admissible enriched baseline at DM p<0.05 (of 8)
\newcommand{\OpsOverEnrichedAdmConc}{-0.19}      % cross-area Spearman of operational improvement vs PV/peak over admissible baseline (broad, n.s.)
\newcommand{\ProxyOverEnrichedAdm}{0.031}        % ERA5 proxy mean MAE improvement over the admissible enriched baseline (information-poor)
% Temporal robustness FY2022-2025 (Phase 33, P6): the cross-area PV->floor gradient
% (H2 first stage) holds in EVERY fiscal year incl. the LNG-crisis FY2022 -> the
% mechanism is structurally stable across regimes, not a single-period artefact.
% (Floor LEVELS are crisis-confounded, so no temporal-growth claim; regime-robustness only.)
\newcommand{\TemporalFloorRhoCrisis}{+0.71}      % FY2022 (LNG-crisis) cross-area Spearman(floor, PV/peak)
\newcommand{\TemporalFloorRhoCrisisP}{0.032}
\newcommand{\TemporalFloorRhoCalm}{+0.83}        % FY2025 (calm) cross-area Spearman(floor, PV/peak)
\newcommand{\TemporalFloorRhoCalmP}{0.005}
\newcommand{\TemporalFloorRhoMin}{+0.58}         % weakest of the four annual gradients (FY2023); all four positive
\newcommand{\KyushuWxValueCrisis}{0.055}         % Kyushu mean weather value, crisis regime FY2022-23 (JPY/kWh)
\newcommand{\KyushuWxValueCalm}{0.038}           % Kyushu mean weather value, calm regime FY2024-25 (JPY/kWh)
% Model-class robustness of the skill-scarcity pivot (Phase 27.5): same
% proxy-vs-operational head-to-head re-run with LightGBM (leaf-wise GBDT).
\newcommand{\LgbmProxyImp}{0.038}                % LightGBM ERA5 proxy mean MAE improvement
\newcommand{\LgbmOpsImp}{0.251}                  % LightGBM operational mean MAE improvement
\newcommand{\LgbmProxyConc}{\rho = +0.48}        % LightGBM proxy cross-area gradient vs PV/peak (n.s.)
\newcommand{\LgbmOpsConc}{\rho = -0.17}          % LightGBM operational cross-area gradient (washout, n.s.)
\newcommand{\CatProxyImp}{0.032}                 % CatBoost ERA5 proxy mean MAE improvement
\newcommand{\CatOpsImp}{0.206}                   % CatBoost operational mean MAE improvement
\newcommand{\CatProxyConc}{\rho = +0.69}         % CatBoost proxy cross-area gradient vs PV/peak (p=0.058, n.s.)
\newcommand{\CatOpsConc}{\rho = +0.26}           % CatBoost operational cross-area gradient (washout, n.s.)
% Threshold-crossing region (Phase 27.0): in curtailment-active (actual
% crossing) slots the operational within-day intensity is significantly
% negative (endogenous selection -> corroborative, not causal).
\newcommand{\CurtailActiveBeta}{-3.47}           % operational within-day intensity, curtail-active slots
\newcommand{\CurtailActiveP}{p = 0.00042}   % asymptotic cluster-t; the wild-cluster bootstrap gives 0.000
\newcommand{\SpearmanPVpeak}{-0.800}             % cross-area PV/peak vs weather Delta-MAE (AIC-best antecedent)
\newcommand{\ProxyFloorFone}{0.087}              % ERA5 proxy floor-event F1 (proxy adds ~no decision value)
\newcommand{\ProxyFloorPRAUC}{0.41}              % ERA5 proxy floor-event PR-AUC
\newcommand{\ProxyFloorBrier}{0.024}             % ERA5 proxy floor-event Brier (cross-fitted)
\newcommand{\ProxyFloorFoneCF}{0.46}             % ERA5 proxy floor F1 at cross-fitted threshold
\newcommand{\OracleFloorBrier}{0.020}            % realised-weather oracle floor-event Brier
\newcommand{\Rsq}{R^2 = 0.537}                  % multivariate FLOOR+PV+Wind (daily)
\newcommand{\Rsqweekly}{R^2 = 0.766}            % multivariate (weekly Phase 22, legacy)
\newcommand{\Rsqfloor}{R^2 = 0.498}             % univariate FLOOR (daily)
\newcommand{\Rsqpv}{R^2 = 0.515}                % univariate PV/peak (daily, AIC-best)
\newcommand{\PanelBeta}{-0.0122}                % per FLOOR-pp, daily n=1260
\newcommand{\PanelAsymCI}{[-0.016, -0.008]}     % asym clustered 90% CI
\newcommand{\PanelBootCI}{[-0.030, +0.001]}     % cluster+block bootstrap 90% CI
\newcommand{\PanelBootCIninetyFive}{[-0.036, +0.005]}  % cluster+block bootstrap 95% CI (Round-20 transparency)

% ----- Phase 24.11.8 noise-degradation sensitivity -----
% Cross-area Spearman attenuates under noise (n=9 loses power) but
% the area-fold panel β with n=1260 clustered SE excludes zero at
% both realistic and 2x-stress noise levels.
\newcommand{\NoiselessSpearman}{-0.733}
\newcommand{\NoiselessPermP}{0.034}
\newcommand{\NoiselessBeta}{-0.0122}
\newcommand{\NoiselessBetaCI}{[-0.016, -0.008]}
\newcommand{\RealisticSpearman}{-0.450}
\newcommand{\RealisticPermP}{0.231}
\newcommand{\RealisticBeta}{-0.0038}
\newcommand{\RealisticBetaCI}{[-0.007, -0.001]}
\newcommand{\StressSpearman}{-0.267}
\newcommand{\StressPermP}{0.490}
\newcommand{\StressBeta}{-0.0050}
\newcommand{\StressBetaCI}{[-0.008, -0.003]}

% ----- Phase 24.14 lag-96 bid-close-admissibility robustness -----
% Pre-shifted weather +48h so harness shift(48) effectively picks t-96h
% weather. Sign preserved at both area-level rank correlation and
% panel beta; panel 90% CI barely excludes zero.
\newcommand{\BcadmSpearman}{-0.300}
\newcommand{\BcadmPermP}{0.438}
\newcommand{\BcadmBeta}{-0.0066}
\newcommand{\BcadmBetaCI}{[-0.013, -0.000]}

% ----- regime labels -----
\newcommand{\FLOOR}{\textsc{floor}\xspace}
\newcommand{\NORMAL}{\textsc{normal}\xspace}
\newcommand{\SPIKE}{\textsc{spike}\xspace}

% ----- DM test -----
\newcommand{\DM}{Diebold--Mariano\xspace}
\newcommand{\dmstat}[2]{$t = #1$, $p = #2$}

% ----- floor/spike thresholds -----
\newcommand{\thrFloorDefault}{$0.5\,\jpykwh$}
\newcommand{\thrSpikeDefault}{$50\,\jpykwh$}

% ---- Phase 43: decision-value sensitivity (2026-07-24, AE restructure) ----
\newcommand{\OpsValueRuleMinK}{149}   % min operational value across 8 threshold rules (max-net-value, zero flag cost), thousand JPY/MW/yr
\newcommand{\OpsValueRuleMaxK}{753}   % max across rules (max-net-value at 1 JPY/kWh flag cost), thousand JPY/MW/yr
\newcommand{\OpsFlagCount}{1466}      % slots flagged by the operational reference rule
\newcommand{\PriceOnlyFlagCount}{1996}% slots flagged by the price-only reference rule (more flags -> per-flag costs amplify operational advantage)

% ---- Phase 44: point-metric broadening, zero refits (2026-07-24, AE restructure) ----
\newcommand{\RmaePriceOnly}{0.673}    % pooled rMAE vs naive-168h, price-only (phase26.1 cache, 8 areas)
\newcommand{\RmaeProxy}{0.662}        % + lag-48 ERA5 proxy
\newcommand{\RmaeOps}{0.594}          % + operational reconstruction
\newcommand{\CacheValueMultipleMae}{7.01}  % histfc/proxy improvement multiple in MAE (same cache)
\newcommand{\CacheValueMultipleRmse}{7.13} % histfc/proxy improvement multiple in RMSE (metric-robust)

% ---- Phase 45: EPEX-DE two-tier external replication (2026-07-24, AE restructure) ----
\newcommand{\DETestDays}{762}      % head-to-head test days, 2024-06-15..2026-07-16
\newcommand{\DEOpsImp}{3.19}       % operational (ICON previous-day run) improvement over price-only, EUR/MWh, DM p<1e-12
\newcommand{\DEProxyImp}{0.19}     % lag-24h ERA5 proxy improvement, EUR/MWh (DM n.s., p=0.29)
\newcommand{\DEOpsOverProxy}{17}   % operational/proxy improvement ratio (17.1)
\newcommand{\DEWindGain}{2.24}     % wind-channel operational gain, EUR/MWh (rank 1/4, ~73% of tier delta)
\newcommand{\DENegGradRho}{+0.121} % exploratory: daily Spearman(op improvement, neg-price share), p=0.0008 -- NOT promoted (pre-committed gate)

% ---- Phase 46: fixed-architecture DNN baseline (2026-07-24, AE restructure) ----
\newcommand{\DnnProxyImp}{-0.010}  % DNN pooled improvement, ERA5 proxy (adds nothing / slightly negative)
\newcommand{\DnnOpsImp}{0.098}     % DNN pooled improvement, operational reconstruction
\newcommand{\DnnGpvImp}{0.109}     % DNN pooled improvement, exact-issuance GPV

% ---- Phase 47: tier-resolved probabilistic evaluation (2026-07-24, AE restructure) ----
\newcommand{\CrpsProxyImp}{0.019}      % 9-area pooled CRPS improvement over price-only, proxy
\newcommand{\CrpsOpsImp}{0.114}        % 9-area pooled CRPS improvement, operational reconstruction
\newcommand{\CrpsValueMultiple}{6.07}  % operational/proxy CRPS-improvement multiple (9-area)

% ---- Phase 49: full-period true previous-day-run tier + extended imbalance window (2026-07-24, AE restructure) ----
\newcommand{\PrevRunImp}{0.219}       % pooled improvement of the true previous-day JMA-MSM run, 21 months, 9 areas (JPY/kWh)
\newcommand{\PrevRunProxyImp}{0.027}  % same-window lag-48 proxy improvement (pooled DM p=0.063, n.s.)
\newcommand{\PrevRunTestDays}{202}    % test days per area, 2024-11-17..2026-07-16 (all 12 calendar months covered)
\newcommand{\PrevRunReconGap}{0.015}  % mean |gap| reconstruction vs true previous-day run (JPY/kWh; DM p=0.35, indistinguishable)
\newcommand{\PrevRunAdmImp}{0.136}    % morning-only strictly gate-admissible subset: true-run improvement (JPY/kWh)
\newcommand{\PrevRunAdmProxyImp}{0.018} % same subset, proxy improvement (7.7x contrast, DM p<1e-6)
\newcommand{\PrevRunDecisionExtK}{110} % extended-window (2024-11..2026-07) true-run decision value, thousand JPY/MW/yr (wedge-regime dependent)

% ============================================================
% Gate-closure revision (2026-07): dissemination-correct vintage ladder,
% FIP commitment decision value, NWP-proxy baseline, tier slope tests.
% Supersedes the post-gate-class headline macros above (retired in the
% R4/R5 rewrite; kept until every reference is migrated).
% ============================================================

% ---- Phase 51: vintage ladder, capacity-weighted primary arm, 9 areas ----
% 732 = 735 calendar days (2024-08-01..2026-08-05) less the 3-day outage
% (2026-04-07/08/09). Do NOT copy the 714/2026-07-18 figures from \CwPanelDays
% below: those belong to the DEMOTED single-point panel, and this comment
% carried both of them for months while the value said 732.
\newcommand{\VintageTestDays}{732}      % daily folds, capacity-weighted arm
\newcommand{\VintageTestSlots}{315{,}414}
% Per-area significance for the stale-observation tier. Methods promises Holm-
% and BH-adjusted p-values "where one test is run per area", and this is that
% family. Only the RAW count existed and Section 5.1 quoted it as eight of nine;
% Holm makes it seven (Shikoku 0.0384 -> 0.0768) while BH leaves it at eight.
\newcommand{\StaleObsSigAreasHolm}{seven}
\newcommand{\StaleObsSigAreasBh}{eight}
\newcommand{\StaleObsImp}{0.047}        % lag-48 ERA5 stale-observation proxy, pooled daily-mean improvement
\newcommand{\StaleObsImpLo}{0.030}
\newcommand{\StaleObsImpHi}{0.063}
\newcommand{\PreGateImp}{0.253}         % pre-gate vintage (previous-day-2 run), disseminated before the gate
\newcommand{\PreGateImpLo}{0.208}
\newcommand{\PreGateImpHi}{0.300}
\newcommand{\PostGateImp}{0.256}        % first post-gate vintage (previous-day-1 run, disseminated 11:30 JST)
\newcommand{\ReconLadderImp}{0.218}     % operational reconstruction (post-gate class)
\newcommand{\ExPostImp}{0.241}          % ex-post realised-weather benchmark (not a ceiling)
\newcommand{\PreGateMultiple}{5.39}     % pre-gate / stale-observation improvement multiple
\newcommand{\PreGateMultipleLo}{4.23}
\newcommand{\PreGateMultipleHi}{7.71}
\newcommand{\GateTimingValue}{0.003}    % post-gate minus pre-gate = value the 10:00 gate forgoes
\newcommand{\GateTimingValueLo}{-0.005}
\newcommand{\GateTimingValueHi}{0.010}
\newcommand{\PreGateVsStaleDiff}{0.206} % Gate-1 quantity
\newcommand{\PreGateVsStaleDiffLo}{0.169}
\newcommand{\PreGateVsStaleDiffHi}{0.246}
\newcommand{\PreGateRetention}{99.0}      % derived: \PreGateImp / \PostGateImp (%)

% ---- Phase 52: FIP commitment decision value (prev2-primary; thousand JPY per MW-installed-year) ----

% ---- Phase 54: NWP-proxy-enriched baseline, 610 days, 8 areas ----
\newcommand{\PreGateOverEnrichedAdm}{0.200} % pre-gate vintage over the fullest admissible baseline
\newcommand{\PreGateOverEnrichedAdmLo}{0.166}
\newcommand{\PreGateOverEnrichedAdmHi}{0.235}
\newcommand{\ProxyCarriedValue}{0.191}      % carried by engineered PV/wind/load proxies from the same vintage
\newcommand{\RawResidualOverProxies}{0.009} % residual value of raw channels on top of those proxies
\newcommand{\RawResidualLo}{0.004}
\newcommand{\RawResidualHi}{0.014}
\newcommand{\EnrichGainAdmDaily}{0.020}     % price-only -> admissible enriched baseline
\newcommand{\ProxyCarriedShare}{95}         % derived: \ProxyCarriedValue / \PreGateOverEnrichedAdm (%)

% ---- Phase 55: tier slope-difference tests ----
\newcommand{\TierDrhoPreGate}{-0.32}    % Spearman-rho difference, stale proxy minus pre-gate vintage
\newcommand{\TierDrhoPreGateLo}{-0.90}
\newcommand{\TierDrhoPreGateHi}{0.05}

% ---- Phase 53: exact archived pre-gate runs (D-2 21/12 UTC), corroboration ----
\newcommand{\ExactPreGateDays}{881}      % delivery days extracted from the GRIB2 archive, 2024-03-01..2026-07-29
\newcommand{\ExactPreGateMargin}{90}     % minutes between the 21 UTC run's dissemination and the gate
\newcommand{\ExactVsPreGateTempRmse}{1.03}   % exact pre-gate vs previous-day-2 vintage, 2-m temperature (degC)
\newcommand{\ExactVsPostGateTempRmse}{0.46}  % exact pre-gate vs first post-gate run, same channel

% ---- Phase 62: single-point proxy validation against OCCTO actuals ----
\newcommand{\ProxyPvRLo}{0.59}          % daytime Pearson r, single-point PV proxy vs realised area PV potential
\newcommand{\ProxyPvRHi}{0.71}
\newcommand{\ProxyPvRMedian}{0.63}
\newcommand{\ProxyWindRLo}{0.02}        % same for the wind proxy: essentially uninformative in most areas
\newcommand{\ProxyWindRHi}{0.42}
\newcommand{\ProxyWindRMedian}{0.14}
\newcommand{\ProxyWindWeakAreas}{six}   % areas (of 8) where the wind proxy correlation is below 0.3
\newcommand{\ProxyDemandRLo}{0.37}      % degree-hour proxy vs realised demand
\newcommand{\ProxyDemandRHi}{0.69}
\newcommand{\CapacityTableImpossible}{seven} % area-channels where realised output exceeds nameplate
\newcommand{\CapacityWorstRatio}{4.8}   % Tokyo wind: realised max / nameplate
\newcommand{\TierDrhoRealisedIdx}{-0.37}     % Delta-rho under the vintage-free saturation index
\newcommand{\TierDrhoRealisedIdxPermArea}{0.14} % exact area-level permutation p, same index

% ---- Phase 61: physical/system metrics + compliant-nomination profit effect ----
% CNPE replaces the retired \FipVoi* family: a value of information cannot be
% negative, so the phase52 quantity is renamed. Supplement S2's theoretical VoI
% is unchanged and stays non-negative.

% ---- Phase 60: the settlement wedge as a continuous treatment ----

% Identity decomposition of the profit effect (supp S2): level vs covariance.
% Floor-vs-curtailment risk ratios, min and max across the eight areas with
% published curtailment actuals. Bound to the two extreme rows of the frozen
% phase25.9 crosstab; the check fails if either stops being the extreme.
\newcommand{\CurtailRiskRatioLo}{19}
\newcommand{\CurtailRiskRatioHi}{96}
% ---- Phase 70: the two channels through which a weather tier reaches the
% commitment. The fourth review objected that the price forecast is not used in the
% nomination. It is, through the curtailment call, and its own contribution is
% POSITIVE -- the whole negative sign is the generation-forecast leg. The two legs
% close to the yen against the reported effect.
% ---- Phase 68: two published prices for the same half-hour. The hypothesis that
% the settlement price K exceeds the system's marginal balancing cost was REFUTED
% (\KAboveCostShare\% of slots overall, none on floor states), so the reported
% inconsistency is between the DAY-AHEAD price and the marginal cost instead.
% Every state is defined by an operational instruction and an exchange outcome,
% never by comparing the two prices under comparison.
% canonical n_boot=10,000. Deterministic given (seed, n_boot) but n_boot-dependent
% (220/254/265 at 50/10,000/2,000 draws), so the manuscript says "more than eight
% months" rather than leaning on the third digit. The verdict -- no date asserted --
% is the same at every draw count.
% ---- Full-sample CNPE (547 test days). The fifth review found the abstract
% headlining the earlier-window -252k while the paper's own position is that no
% break date is identifiable and the wedge is continuous; a paper that argues the
% split is arbitrary cannot headline one side of it. The full sample is primary
% from here on; the windows are sensitivity.
% ---- Phase 69: netting. The reviewer's first objection was that one party's
% gross deviation is not system imbalance. With the wide-area block's own signed
% imbalance volume (phase71) our reduction can be split by the direction of the
% system's own error. Only the NET figure may be reported as a headline; the gross
% figure must never appear unqualified again.
% ---- Phase 74: capacity-weighted multi-point re-estimation (fifth review #2) ----
% Gate 5 reproduced phase59's nine cells EXACTLY (max|diff| = 0 over the full
% panel), Gate 35 fired (CW proxy wins 7 of 8 areas), and Gate 36 landed on the
% pre-registered 7b branch: the gate-timing null holds under the better
% measurement and one-tick equivalence is established there too. The measured
% interaction between measurement method and lead -- the reviewer's conjecture --
% is indistinguishable from zero.
% The headline tier's lead range at valid time. A deterministic property of the
% run schedule (12 UTC on D-2 = 21:00 JST, delivery-day slots 00:00-23:30 JST),
% not a measurement, so it is an exempt protocol constant. Defined once and used
% in BOTH the admissibility table and Section 7, which used to hand-type the
% range in the table and describe the tier as "about 48 hours old" in the text.
% The clock time the excluded run reaches users, on the published schedule.
% Needed because Section 5.3 now argues from the ORDER of two events: the
% auction clears at 10:00 and this run arrives after it, so no price in the
% sample can embody it.
\newcommand{\PostGateArrival}{11:30}

% ---- Phase 96: the same contrast on LEAR instead of XGBoost ----
% Section 4.4 named LEAR as a scale anchor and it had never run on the 732-fold
% panel, so a reader could not tell whether the headline increment is a property
% of the information set or of the learner. It is both: direction and breadth
% replicate (9/9 areas, p<1e-4) while the magnitude halves.
\newcommand{\LearBaseMae}{2.216}
\newcommand{\LearPreGateMae}{2.088}
\newcommand{\LearImp}{0.128}
\newcommand{\LearImpLo}{0.090}
\newcommand{\LearImpHi}{0.168}
\newcommand{\LearAreasImproved}{nine}

% The dependence-aware standard error of the gate-timing contrast. Printed so
% Section 5.3 can say plainly that the equivalence margin and the design's own
% resolution are the same order of magnitude -- a referee will compute this.
\newcommand{\GateTimingSe}{0.0038}

% ---- Phase 99: audit rates on the full-text subset ----
% A referee objected that the denominators partly measure TEXT ACCESS rather than
% practice. On the subset whose full text we read, the gate-time rate RISES while
% publication time and admissibility checking fall to zero -- reading everything
% confirms the absence instead of finding the disclosures.

% ---- Phase 98: the dissemination audit in full, and the poll-grid test ----
% A referee offered a fixed polling cadence as an equally consistent explanation
% of the mirror timestamps. It is testable and it fails: the lags vary at second
% resolution inside a band a minute wide, which a poll schedule cannot produce.
\newcommand{\DissemDistinctLags}{66}
\newcommand{\DissemDistinctFracMin}{52}
\newcommand{\DissemGapSeconds}{1}
\newcommand{\DissemWithinMinutePct}{99.86}

% ---- The gate-timing contrast on the DAYLIGHT window, primary arm ----
% Three quantities in this article sit within a thousandth of each other -- the
% product contribution (0.013), the ex-post-below-pre-gate gap (0.012), and this
% one -- and the article called the first "small but detectable" while calling the
% gate-timing effect indistinguishable from zero. Both statements are true, and
% together they read as a contradiction unless the article says why: the difference
% is PRECISION, not magnitude. The product contribution's interval excludes zero on
% the full panel; this one, on a third of the slots, does not. Its point estimate
% also exceeds one tick, which the equivalence sentence has to acknowledge, so it is
% bound and printed rather than left in the artefact.
\newcommand{\GateTimingDaylight}{0.013}
\newcommand{\GateTimingDaylightLo}{-0.006}
\newcommand{\GateTimingDaylightHi}{0.032}
\newcommand{\GateTimingDaylightP}{0.19}

% ---- Phase 108: the cost--loss increment, against a fallback anyone would use ----
% The article reported the weather channels' contribution as V(with) - V(without)
% and headlined the cost--loss ratio where that difference is largest (0.64). Above
% alpha = 0.44 the without-weather forecast scores NEGATIVE -- worse than always or
% never protecting -- so no actor would run it; they would fall back to climatology,
% which scores zero. Differencing against a dominated alternative credits the weather
% channels with that alternative's losses, and the headlined alpha is exactly where
% the credit is largest. 28 of the 49 grid points were dominated.
%
% Worse, the article read the curve backwards. V(with) falls 22-fold across the grid,
% so the vintage is worth MOST to actors whose protection is cheap relative to their
% exposure; the article concluded it serves actors with expensive protection.
\newcommand{\CostLossDominatedAlpha}{0.44}   % where the no-weather score goes negative
\newcommand{\CostLossPeakAlphaFixed}{0.42}   % peak increment against a USABLE fallback
\newcommand{\CostLossPeakGainFixed}{0.319}
\newcommand{\CostLossVLow}{0.76}             % value score at the cheapest protection
\newcommand{\CostLossVHigh}{0.03}            % and at the dearest

% ---- Phase 106: the two-target asymmetry in COMMENSURABLE units ----
% The article led on a pair of variance shares -- 14.0% of the physical residual's
% variance against 1.2% of the price residual's -- and three independent referees
% objected that those are shares of two DIFFERENT dependent variables with very
% different residual variances, so their ratio means nothing. They were right, and
% they understated the fix: comparing each target's own fractional error reduction
% is both commensurable and about four times the contrast. The variance shares stay
% where they belong, inside the mechanism regression on each target separately;
% they no longer carry the headline.
\newcommand{\PriceRelReduction}{0.14}   % gate-timing gain as a share of the price arm's own MAE (%)
\newcommand{\AsymmetryRatio}{50}        % physical over price, in commensurable units

% ---- Phase 102: signal or learner, and the window mismatch it exposed ----
% The pre-registered question was whether the post-gate price null is a property of
% the SIGNAL or of the LEARNER. Handed the run-to-run innovation directly as four
% features, the same learner recovers 99.3% of the physical gain (positive control)
% and still moves pooled price MAE by nothing detectable -- so the null is the
% signal's, not the estimator's.
%
% What the same run exposed is more important, and it is a defensive necessity
% rather than a bonus: the article compares a DAYLIGHT-ONLY physical number against
% an ALL-SLOT price number. On the physical claim's own window the price gain is
% significant. The restriction decomposes cleanly -- the HOURS carry it, the AREAS
% do not -- so the pooled null is in part a dilution over 34 night slots in which
% an update to the cloud field cannot move photovoltaic-driven price formation.
% "No price-relevant content" is therefore too strong; the defensible claim is
% "none detectable in the 24-hour pooled forecast".
\newcommand{\InnovPriceDelta}{0.003}      % innovation as features, pooled 24-hour price gain
\newcommand{\InnovPriceLo}{-0.002}
\newcommand{\InnovPriceHi}{0.008}
\newcommand{\InnovPriceP}{0.216}          % HLN-DM, L=3
\newcommand{\InnovNoiseCost}{-0.003}
\newcommand{\InnovPhysRecoveredPct}{99.3}  % EXEMPT-worthy: share of the physical gain the innovation recovers      % four matched-variance noise columns: the dimensionality yardstick
\newcommand{\DaylightPriceDelta}{0.015}   % same contrast on the physical claim's OWN window (09-15 JST)
\newcommand{\DaylightPriceLo}{0.004}
\newcommand{\DaylightPriceHi}{0.026}
\newcommand{\DaylightPriceP}{0.010}       % HLN-DM, L=3
\newcommand{\HoursOnlyDelta}{0.015}       % restriction decomposition: hours carry it
\newcommand{\AreasOnlyDelta}{0.004}       % areas do not
\newcommand{\AreasOnlyLo}{-0.0005}
\newcommand{\AreasOnlyHi}{0.009}

% ---- Phase 103b: the floor channel carries a third, not most ----
% Claim three said weather value concentrates in floor transitions. Decomposed
% exactly (max residual 4e-16) into a hurdle system, the floor-STATE channel
% carries 31.7% and the conditional-LEVEL channel 70.7%. The strong aggregate
% reading is not supported; the article's actual sentence -- an intensity on a
% 3.33% slice -- is exactly consistent, and 3.33% x 7.35 ~= 24% is the arithmetic
% that sentence had left unstated. The floor classifier itself improves decisively
% (Brier 0.0246 -> 0.0173), so the modest share is prevalence, not a weak signal.
\newcommand{\HurdleFloorShare}{31.7}      % floor-state channel, % of the pooled gain
\newcommand{\HurdleFloorShareLo}{26.3}
\newcommand{\HurdleFloorShareHi}{37.1}
\newcommand{\HurdleLevelShare}{70.7}      % conditional-level channel, % -- the majority conduit
\newcommand{\HurdleFloorChannel}{0.080}   % the floor-state channel in JPY/kWh
\newcommand{\HurdleLevelChannel}{0.179}
\newcommand{\HurdleBrierBase}{0.0246}     % floor classifier without weather
\newcommand{\HurdleBrierWx}{0.0173}       % with it
\newcommand{\HurdleResidual}{-2.4}        % the THIRD term. Dropping it made the printed shares sum to 102.4% under a sentence claiming they sum exactly
\newcommand{\HurdleKyushuShare}{55.9}     % the only area where the floor channel is the majority

% ---- Phase 101: the physical contrast in normalised units ----
% Methods called a capacity-weighted MW mean a "scale-free comparison". It is not:
% weighting by capacity makes the mean a statement about the largest fleets. Worse,
% read literally the phrase names the arm that is 17% LARGER, so the loose wording
% favoured the paper. These give both weightings explicitly and let the text name
% which one it means.
%
% The STEP is what survives: across three registers the area-equal step spans
% 0.70-0.75 pp and the capacity-weighted 0.83-0.87 pp. The LEVEL does not survive
% that well (it moves about 1.1 pp), so the level is quoted only as an honest
% statement about a coarse instrument, never as a normalised result.
\newcommand{\PhysNmaeStepEq}{0.71}      % nMAE step, areas weighted EQUALLY (genuinely scale-free)
\newcommand{\PhysNmaeStepEqLo}{0.58}
\newcommand{\PhysNmaeStepEqHi}{0.85}
\newcommand{\PhysNmaeStepCw}{0.83}      % nMAE step, CAPACITY-weighted -- the published MW figure over mean capacity
\newcommand{\PhysNmaeStepOccto}{0.70}   % same step on the OCCTO register
\newcommand{\PhysNmaeStepMaxreal}{0.75} % same step on a vintage-free denominator
\newcommand{\PhysNmaeLevelEq}{12.3}     % pre-gate LEVEL, areas equal: a coarse instrument, not a PV forecast
\newcommand{\PhysRelEq}{5.8}            % relative improvement, areas weighted equally (%)
\newcommand{\PhysRelEqLo}{4.8}
\newcommand{\PhysRelEqHi}{6.9}
\newcommand{\PhysRelPooled}{6.8}        % relative improvement, pooled ratio of means (%)
\newcommand{\PhysRelAreaMin}{2.7}       % weakest area (Hokkaido)
\newcommand{\PhysRelAreaMax}{8.1}       % strongest area (Kyushu)

% ---- Phase 55: the cross-area gradient, which stays UNDETERMINED ----
% These four were LITERALS in the Results text, and one of them was wrong: the
% prose said the forecast tiers correlate "between +0.40 and +0.53", but +0.53 is
% the EX-POST realised-weather tier, which the article is careful to separate from
% the forecast tiers everywhere else. The forecast tiers span +0.40 to +0.43. The
% inflated range made a null result look more suggestive than it is.
\newcommand{\CrossAreaRhoProxy}{0.12}        % stale-observation proxy (era5)
\newcommand{\CrossAreaRhoForecastLo}{0.40}   % lowest forecast tier (histfc)
\newcommand{\CrossAreaRhoForecastHi}{0.43}   % highest forecast tier (prev2)
\newcommand{\CrossAreaRhoExPost}{0.53}       % ex-post tier -- NOT a forecast tier
\newcommand{\CrossAreaMinP}{0.14}            % smallest p in the family; supports "none below 0.05"

% ---- Phase 105: the excluded run's arrival, file by file, against its own gate ----
% The referee's first-priority objection was that the paper judged the D-1 00 UTC
% run inadmissible on a MEDIAN lag, and that a median says nothing about the tail:
% perhaps some forecast-hour files, on some days, did clear 10:00 JST. Asked at
% file resolution against each delivery day's own gate, the answer is that none
% ever did -- and the decisive figure is not the count but the SPREAD. Over 2.5
% years the earliest and latest of 2,643 files differ by under two minutes, so the
% exclusion is structural rather than probabilistic and no tail argument survives.
\newcommand{\GateAuditFiles}{2{,}643}       % D-1 00Z files audited: 881 days x 3 horizon files
\newcommand{\GateFastestFileMin}{75.6}      % earliest arrival in the whole archive, minutes AFTER the gate
\newcommand{\GateSlowestFileMin}{77.3}      % latest arrival, minutes after the gate
\newcommand{\GateArrivalSpreadMin}{1.7}     % total spread across 2.5 years -- the number that ends the tail argument
\newcommand{\GateFastestArrivalJst}{11:15:35}  % clock time of that earliest file
\newcommand{\GateFilesPreGate}{zero}        % files that beat the gate, out of \GateAuditFiles
\newcommand{\GateDaysAnyPreGate}{zero}      % days on which ANY file beat the gate, out of 881
\newcommand{\GateFhSpreadSec}{15}           % spread of per-horizon-file MEDIANS: refutes an early-arriving subset

% ---- Phase 100: which cycles the mirror carries beyond FH39 ----
% The archive-gap caveat rested on a code comment and the memory of a manual
% probe; no artefact recorded a 404 and no code had ever requested 18Z x FH40-51.
% 280 HEAD requests over 8 cycles x 5 horizon files x 7 dates settle it, with
% 00Z x FH40-51 as a positive control so a broken probe cannot read as a finding.
\newcommand{\ProbeRequests}{280}            % HEAD requests issued
\newcommand{\ProbeCycles}{eight}            % initialisation cycles probed (00,03,...,21 UTC)
\newcommand{\ProbeLongHorizonCycles}{two}   % cycles carrying FH40-51 and FH52-78: 00 and 12 UTC only
\newcommand{\ProbeDates}{seven}             % delivery dates spanning the window and the model upgrade

% ---- Phase 97: decision value of the floor forecasts (Murphy cost-loss) ----
% Section 3 invokes the cost-loss tradition and the article never ran one. The
% value score is 1 for a perfect forecast and 0 for one no better than always or
% never protecting; alpha = C/L indexes the actor. Without pre-gate weather the
% forecast is only useful to actors with cheap protection; with it, every ratio
% in the grid gains.
\newcommand{\CostLossMaxV}{0.76}
\newcommand{\CostLossNoWxCutoff}{0.42}
\newcommand{\CostLossPeakAlpha}{0.64}
\newcommand{\CostLossPeakGain}{0.36}
\newcommand{\CostLossPeakGainLo}{0.26}
\newcommand{\CostLossPeakGainHi}{0.48}
\newcommand{\CostLossBaseRate}{3.4}

% ---- Phase 95: what "nine of nine areas" actually rests on ----
% Interconnection capacity is allocated through the spot auction, so area prices
% coincide when no interface binds. The counts the paper leans on are about the
% per-area IMPROVEMENT series, whose effective independence is what matters.
\newcommand{\AreaUnsplitPct}{8.2}
\newcommand{\AreaImpMeanCorr}{0.338}
\newcommand{\AreaEffN}{4.11}
\newcommand{\AreaEffNPrices}{1.53}

% The D-2 18 UTC cycle: a fresher whole-day pre-gate rung than the headline one,
% absent from the ladder because THIS ARCHIVE stops the intermediate cycles at
% forecast hour 39. Section 4.2 used to call that a fact about the observing
% system; it is a fact about the mirror, and the cycle curve already scores this
% cycle on the common daytime window.
\newcommand{\CycleEighteenZMarginHrs}{4.73}
\newcommand{\CycleEighteenZImp}{0.477}
\newcommand{\CycleEighteenZImpLo}{0.375}
\newcommand{\CycleEighteenZImpHi}{0.587}

% ---- Phase 94: the price regime the evaluation window actually is ----
% Section 2.1 cited the FY2022 LNG crisis, which is TRAINING data only, and so
% implied spike coverage the evaluation era does not have. Near-floor slots
% outnumber scarcity slots by about 99 to 1 here, so every headline number is a
% floor-regime result and the paper now says so.
\newcommand{\EvalPNinetyFive}{20.7}
\newcommand{\EvalPNinetyNine}{30.2}
\newcommand{\EvalNearFloorPct}{3.33}
\newcommand{\EvalScarcityPct}{0.034}
\newcommand{\EvalFloorScarcityRatio}{99}

\newcommand{\PreGateLeadLo}{27}
\newcommand{\PreGateLeadHi}{50}

\newcommand{\CwPoints}{46}             % prefecture points across nine areas
\newcommand{\CwProxyWins}{7}           % areas where the CW proxy beats single-point
\newcommand{\CwProxyAreas}{8}
\newcommand{\CwProxyKyushuSingle}{0.883}
\newcommand{\CwProxyKyushuCw}{0.916}
% Enriched-baseline lineage branches (phase78/phase54 variants)
\newcommand{\EnrichedDays}{698}
\newcommand{\EnrichedDaysKyushu}{668}  % the one area with a shorter enriched panel
\newcommand{\PreGateOverEnrichedDrop}{0.205}
\newcommand{\PreGateOverEnrichedLagX}{0.197}
% Robustness block (phase80/82/83/78)
\newcommand{\VersionInteraction}{-0.006}
\newcommand{\VersionInteractionLo}{-0.022}
\newcommand{\VersionInteractionHi}{0.011}
% Pre-gate improvement either side of the 2025-05-20 MSM physics upgrade
\newcommand{\VersionPreV}{0.302}
\newcommand{\VersionPostV}{0.220}
\newcommand{\VersionPreVDays}{292}
\newcommand{\VersionPostVDays}{440}
\newcommand{\VersionPreGateDelta}{-0.082}
\newcommand{\VersionPreGateDeltaLo}{-0.178}
\newcommand{\VersionPreGateDeltaHi}{0.016}
\newcommand{\VersionPreGateDeltaP}{0.095}
\newcommand{\VersionLagMedianV}{135.8}
\newcommand{\VersionLagMedianPostV}{135.9}
\newcommand{\HeadlineDemandW}{0.266}
\newcommand{\HeadlinePvW}{0.277}
\newcommand{\PolitisWhiteBlock}{2.0}
\newcommand{\CompositeWorstMarginMin}{103.6}
% Expectations channel (phase79): physical vs price target
% RETIRED 2026-08-13: phase79's physical-target numbers are superseded by
% phase89, which cross-fits the scale (phase79 fitted it on the full sample),
% restricts to the panel the price ladder shares, and corrects the OCCTO stamp
% convention per operator. The replacements are \PhysDelta and the \Phys*
% block below. Kept as a comment rather than deleted so a reader of the git
% history can see what the numbers were and why they moved.
% Single-point sensitivity arm (phase85 arm_sp_robustness; CW is primary)
\newcommand{\SpPreGateImp}{0.225}
\newcommand{\SpPreGateImpLo}{0.183}
\newcommand{\SpPreGateImpHi}{0.269}
\newcommand{\SpGateTiming}{0.004}
\newcommand{\SpGateTimingLo}{-0.003}
\newcommand{\SpGateTimingHi}{0.010}
\newcommand{\CwUplift}{0.028}          % CW minus single-point, headline pre-gate tier
\newcommand{\CwUpliftPost}{0.027}      % same for the post-gate tier: near-identical
\newcommand{\CwInteraction}{-0.0009}   % (CW gate step) - (single-point gate step)
\newcommand{\CwInteractionLo}{-0.0069}
\newcommand{\CwInteractionHi}{0.0052}
\newcommand{\CwInteractionP}{0.77}
\newcommand{\CwPanelDays}{714}
% ---- Phase 73: cycle-by-cycle curve on the common 08:00-16:00 window ----
\newcommand{\CycleDays}{735}          % common-window panel days
\newcommand{\CycleStepUp}{0.072}      % D-3 12Z -> D-2 12Z improvement step, JPY/kWh
\newcommand{\CycleStepUpLo}{0.048}
\newcommand{\CycleStepUpHi}{0.098}
\newcommand{\CycleFlatMaxStep}{0.008} % largest adjacent step after D-2 12Z (all n.s., < 1 tick)
\newcommand{\CycleSlope}{-0.0024}     % Gate 41 linear slope on gate-margin hours (41b)
\newcommand{\CycleSlopeLo}{-0.0032}
\newcommand{\CycleSlopeHi}{-0.0016}
\newcommand{\CycleGateStepCommon}{0.006}  % common-window gate step
\newcommand{\CycleGateStepCommonLo}{-0.009}
\newcommand{\CycleGateStepCommonHi}{0.022}
\newcommand{\CycleStepsAfterFirst}{3}   % adjacent steps after the D-3 -> D-2 12Z rise
\newcommand{\CycleStepsEquiv}{1}        % of those, how many establish one-tick equivalence
\newcommand{\CycleWidestCiTicks}{1.7}   % widest post-first-step interval, in ticks
\newcommand{\CycleGateStepTostP}{0.019} % one-tick TOST on the gate step itself (bootstrap arm)
\newcommand{\CycleGateStepTostPDm}{0.022}   % the HLN-DM L=3 arm of the same test
% The two-prices ratio, printed rather than rounded to "orders of magnitude":
% 9.18 / 0.020 = 459x, which is 2.66 orders, not three.
\newcommand{\CyclePlateauMarginHrs}{10.7} % measured dissemination-to-gate margin where the plateau begins
% ---- Phase 72: dissemination timing measured from archive timestamps ----
% ---- Phase 92: photovoltaic curtailment days per fiscal year ----
% Section 2 said "roughly 30 curtailment days per year" for Kyushu, inherited
% verbatim from a superseded draft where it described an FY2018-era level. On
% the records this paper scores against it is about five times that, and the
% gap to the next area is a factor of two rather than "far more".
\newcommand{\CurtailDaysKyushuFyOne}{148}  % FY2024
\newcommand{\CurtailDaysKyushuFyTwo}{156}  % FY2025
\newcommand{\CurtailDaysKyushuRound}{150}  % the two-year mean, rounded for prose
\newcommand{\CurtailDaysRunnerUp}{81}      % next most affected area, two-year mean
\newcommand{\CurtailRatioTopNext}{1.9}     % Kyushu / next area

% The PUBLISHED schedule, against which the measured lag is 14 minutes shorter.
% Bound because Section 2.2 now leans on the gap between the two: the schedule is
% not a floor, so it cannot by itself prove the excluded run was late.
\newcommand{\DissemSpecLagMin}{150}
\newcommand{\DissemMedianLagMin}{136}   % median per-run completion lag, minutes after initialisation (= 2 h 16 m)
% The h/min rendering of the line above. It was hand-typed in the abstract and
% in Section 2 while the minute figure next to it was bound, so a re-measured
% lag would have moved one and left the other. Checked by the guard's hhmm mode
% against the same provenance field, so the two renderings cannot diverge.
\newcommand{\DissemMedianLagHM}{2~h~16~min}
\newcommand{\DissemMedianLagHMJa}{2 時間 16 分}
\newcommand{\DissemMaxLagMin}{159}      % worst single day (12Z cycle), minutes
\newcommand{\DissemAuditDays}{881}      % delivery days audited
\newcommand{\DissemAuditFiles}{7048}    % archived files probed, none missing
\newcommand{\DissemWorstMarginTwentyOneZHrs}{1.7} % worst pre-gate margin, D-2 21Z (gate - init = 4 h)
\newcommand{\DissemWorstMarginTwelveZHrs}{10.3}    % worst pre-gate margin, D-2 12Z (gate - init = 13 h)
% ---- Phase 75: decision-scale equivalence (fifth review #1 landing) ----
% ---- Phase 66: floor-event probabilistic evaluation. The conceptual framework
% says weather predicts a REGIME, not only a level; until this phase that claim
% was only ever evaluated by a pooled MAE. Threshold 0.5 JPY/kWh, all hours,
% macro-averaged across the nine areas (slot-pooling makes the PR curve a
% base-rate mixture dominated by Kyushu and Shikoku).
\newcommand{\FloorPrAucBase}{0.339}      % price-only floor PR-AUC, macro-avg
\newcommand{\FloorPrAucPreGate}{0.635}   % pre-gate run
\newcommand{\FloorPrAucOracle}{0.622}    % realised-weather oracle: BELOW pre-gate
% The floor contrast on BOTH measurements. The primary (capacity-weighted)
% values carry the paper; the single-point ones are the sensitivity arm, and
% they are what an earlier version printed beside capacity-weighted LEVELS.
\newcommand{\FloorPrAucDelta}{0.296}     % pre-gate minus price-only
\newcommand{\FloorPrAucDeltaSp}{0.276}
\newcommand{\FloorPrAucDeltaSpLo}{0.190}
\newcommand{\FloorPrAucDeltaSpHi}{0.345}
\newcommand{\GateTimingPrAucSp}{0.008}
\newcommand{\GateTimingPrAucSpLo}{-0.026}
\newcommand{\GateTimingPrAucSpHi}{0.045}
\newcommand{\FloorPrAucDeltaLo}{0.243}
\newcommand{\FloorPrAucDeltaHi}{0.336}
\newcommand{\FloorPrAucDeltaEndpointLo}{0.205}
\newcommand{\FloorPrAucDeltaEndpointHi}{0.367}
\newcommand{\FloorPrAucAreasPos}{nine}   % areas where the contrast is positive
\newcommand{\FloorBrierDelta}{-0.0072}   % Brier improvement, DM-testable
\newcommand{\FloorLiftTokyo}{57.2}       % base-rate lift: LOWEST level, HIGHEST lift
\newcommand{\FloorLiftKyushu}{12.4}
\newcommand{\FloorRecallKyushuBase}{0.186}  % recall at FPR <= 0.01, price-only
\newcommand{\FloorRecallKyushuOp}{0.722}    % same, pre-gate run
\newcommand{\GateTimingPrAuc}{0.010}     % gate timing in the TAIL: also a null
\newcommand{\GateTimingPrAucLo}{-0.002}
\newcommand{\GateTimingPrAucHi}{0.023}
% REBOUND 2026-08-14 to the CAPACITY-WEIGHTED midday scope, which the full grid
% finished computing at 01:12. Until then this was the only floor figure in the
% paper still measured on the demoted single-point arm, and the text said so.
\newcommand{\FloorMiddayDelta}{0.298}    % midday contrast, capacity-weighted
\newcommand{\FloorMiddayDeltaSp}{0.279}  % the single-point arm, kept as sensitivity
% ---- Phase 67: equivalence tests (TOST). Margin = one JEPX price tick, which
% every realised price in the panel lies on exactly. Applied to all four nulls,
% including the two where it FAILS -- "not distinguishable from zero" is not the
% same claim as "shown equivalent to zero".
\newcommand{\GateTimingEquivMargin}{0.01}   % one JEPX price tick, JPY/kWh
% The same margin as a share of the price-only baseline MAE (phase83 / phase82).
\newcommand{\GateTimingMarginPctBaseline}{0.47}
\newcommand{\GateTimingTostP}{0.024}        % TOST p, HLN-DM L=3, tick margin
\newcommand{\GateTimingTostPBoot}{0.023}    % same by the paired block bootstrap
\newcommand{\GateTimingCiNinetyLo}{-0.004}  % 90% interval (the TOST interval)
\newcommand{\GateTimingCiNinetyHi}{0.009}
\newcommand{\GateTimingMarginUtil}{86}      % 90% bound as a % of one tick: tight
\newcommand{\CompositeTostP}{0.751}         % NOT equivalent: exceeds one tick
\newcommand{\IngestionTostP}{0.0000001}     % equivalent, comfortably
\newcommand{\GateTimingAreasEquiv}{zero}    % areas equivalent at one tick alone

% ---- magnitude twins for signed quantities -------------------------------
% The signed macros are correct against provenance, but prose like "cuts X by
% \CnpeAbsImbPre" renders a double negative. These carry the magnitude so the
% verb can carry the sign; the CI endpoints swap order when negated.

% ---- Phase 59: run-pinned ladder (replaces the composite-tier figures above) ----
% The headline pre-gate tier is now the D-2 12 UTC run, pinned to one
% initialisation. Gate 7 landed on the pre-registered 7b branch: the run the gate
% excludes is not distinguishable from zero, so it is reported as an upper bound.
\newcommand{\HeadlineTierName}{the $D{-}2$ 12\,UTC run}
\newcommand{\GateCompositeImp}{0.254}       % per-slot freshest admissible run (corroboration)
\newcommand{\GateCompositeImpLo}{0.209}
\newcommand{\GateCompositeImpHi}{0.301}
\newcommand{\GateTimingUpperBound}{0.010}   % 95% upper bound on what a later gate would buy
\newcommand{\GateTimingP}{0.50}             % pooled DM p, L=3 (n.s.; L=9 gives 0.40)
\newcommand{\VintageEffectWithinEraFive}{0.194}  % product held fixed: realised minus lag-48 ERA5
\newcommand{\VintageEffectWithinEraFiveLo}{0.159}
\newcommand{\VintageEffectWithinEraFiveHi}{0.230}
\newcommand{\ProductBound}{0.013}           % vintage ~fixed: run-pinned NWP minus realised ERA5
\newcommand{\ProductBoundLo}{0.002}
\newcommand{\ProductBoundHi}{0.023}
\newcommand{\CompositeVsPinned}{-0.013}     % old composite minus the run-pinned headline
\newcommand{\CompositeVsPinnedLo}{-0.022}
\newcommand{\CompositeVsPinnedHi}{-0.005}
\newcommand{\CompositeVsPinnedP}{0.003}
\newcommand{\IngestionDiff}{0.000}         % same initialisation, pinned run vs composite
\newcommand{\IngestionDiffP}{0.98}
\newcommand{\DoseSlopePerLeadHour}{-0.0002} % improvement per extra hour of lead, run-pinned only
\newcommand{\HeadlinePerAreaLo}{0.092}      % Hokkaido
\newcommand{\HeadlinePerAreaHi}{0.378}      % Tohoku

% ---------------------------------------------------------------------------
% Literature audit of weather-input timing (phase88; supp. S7 + Section 2.3).
% Rates are over the papers whose published text could be read; papers that
% could not be read are reported separately and never counted as evidence.
% ---------------------------------------------------------------------------
\newcommand{\LitScreened}{55}
\newcommand{\LitIncluded}{43}
% Studies whose audit trail records the BODY was read. The published robustness
% check sliced on access_level instead, which records whether the PUBLISHER's
% version was reachable -- so a study whose ScienceDirect page 403'd while its
% accepted arXiv manuscript was read cover to cover counted as not-full-text. On
% that wrong slice the subset was 20 and publication time read as zero; on the
% right one it is 42 of 43 and the rates are the headline rates. A referee found
% it. The corrected reading is the stronger argument: the audit EXCLUDED studies
% it could not read, so there is no unread stratum for the access objection to
% hide in, and the four LitFullText* macros are retired rather than repaired.
\newcommand{\LitReadN}{42}

% ---- The panel's own start date, which the body never stated ----
% Two referees noticed the consequence: the article asserted the FY2022 LNG spike
% window "falls entirely in training data", and it does not -- the panel begins
% 2024-03-02, so FY2022 is outside the sample altogether. The earliest folds also
% train on about four months rather than years, which a reader is entitled to know.
\newcommand{\PanelStart}{2024-03-02}
\newcommand{\LitExcluded}{5}
\newcommand{\LitUnresolved}{7}
\newcommand{\LitFullText}{20}
\newcommand{\LitPartial}{17}
\newcommand{\LitAbstractOnly}{6}
\newcommand{\LitGateTimePct}{42}
\newcommand{\LitRunIdPct}{5}
\newcommand{\LitInitPct}{2}
\newcommand{\LitPubTimePct}{5}
\newcommand{\LitAdmissiblePct}{5}
\newcommand{\LitRevisionPct}{0}
\newcommand{\LitAdmissibleN}{2}
\newcommand{\LitRevisionNotStated}{36}
\newcommand{\LitRecordsCollected}{57}   % before duplicate works were collapsed
\newcommand{\LitRevisionPartial}{6}
\newcommand{\LitAdmissiblePartial}{24}
\newcommand{\LitPubTimePartial}{17}

% ---- Phase 110: per-slot forecast dispersion (Paper A; field subset of phase 75)
% These two are rendered in the main article and are NOT decision quantities:
% they are what pooled-MAE equivalence cannot bound. They used to sit inside the
% paper-B block because their source artefact also carries the commitment gates,
% which made Paper A cite two numbers whose provenance the public package could
% not ship. phase110_slot_dispersion.json carries only these.
\newcommand{\DecShareGtTickPct}{97.7}   % share of slots where the two runs' price forecasts differ > 1 tick
\newcommand{\DecMeanAbsDy}{0.48}        % mean per-slot |yhat_pre - yhat_post|, JPY/kWh

% ===== BEGIN paper-B-only macros ==========================================
% Settlement-paper macros. They stay DEFINED here so the archived Japanese
% reading copy and the preserved settlement sections still compile, and so
% the drift guard keeps verifying them against provenance that has not moved.
% They are not used by the forecast-vintage manuscript. The guard asserts that
% the set of macros between these markers equals the set of checks tagged
% paper_b=True, in BOTH directions, so this block cannot silently drift out of
% step with what the settlement paper actually claims -- and extracting the
% settlement paper later is a matter of moving the block, not of re-deciding
% what belongs to it.
\newcommand{\FloorWedgeMedianPre}{5.94} % median (imbalance price - day-ahead price) on floor slots, pre-collapse
\newcommand{\FloorWedgeMedianPost}{-0.01}
\newcommand{\AlwaysHedgeBeatPerfectK}{19}         % never-sell policy minus perfect information, pre-collapse
\newcommand{\AlwaysHedgeBeatPerfectMatchedK}{330} % same, on the matched committed-window days
\newcommand{\CnpeAbsImbPre}{-127}       % Delta absolute imbalance, MWh per MW-installed-yr, earlier window
\newcommand{\CnpeAbsImbPreLo}{-155}
\newcommand{\CnpeAbsImbPreHi}{-98}
\newcommand{\CnpeAbsImbPrePct}{28.5}    % same, as a percentage reduction
\newcommand{\CnpeAbsImbPost}{-60}       % later window
\newcommand{\CnpeAbsImbPostLo}{-100}
\newcommand{\CnpeAbsImbPostHi}{-21}
\newcommand{\CnpeAbsImbPooled}{-109}    % full window
\newcommand{\ChatMaePreOp}{0.0096}      % curtailment-share forecast MAE, operational information set
\newcommand{\ChatMaePreBase}{0.0132}    % same, price-only
\newcommand{\AlwaysHedgeAbsImbPre}{1{,}289} % never-sell policy absolute imbalance, MWh per MW-yr
\newcommand{\CnpePreK}{-252}            % compliant-nomination profit effect, thousand JPY per MW-yr
\newcommand{\CnpePreAbsK}{252}          % magnitude, for prose that already carries the sign
\newcommand{\CnpePreLoK}{-415}
\newcommand{\CnpePreHiK}{-109}
\newcommand{\CnpePostK}{+48}
\newcommand{\CnpePostLoK}{-84}
\newcommand{\CnpePostHiK}{+167}
\newcommand{\CnpeBandSigEps}{0.05}      % largest nomination band with a significantly negative CNPE
\newcommand{\CnpeBandFlipEps}{0.20}     % band at which the point estimate turns positive
\newcommand{\CnpeAbsImbBandTwenty}{-59} % Delta absolute imbalance at that band: shading degrades physics
\newcommand{\FloorSlotSharePct}{3.0}    % floor slots as a share of all slots
\newcommand{\FloorProfitSharePct}{58.3} % share of the profit difference they carry
\newcommand{\WedgeBetaFloor}{-37.7}     % JPY per MW-day per JPY/kWh of floor-slot wedge exposure
\newcommand{\WedgeRsqFloorPct}{41}      % R^2 of the floor-slot profit component
\newcommand{\WedgeRsqNonfloorPct}{2}    % R^2 of the non-floor component
\newcommand{\WedgePostP}{0.26}          % post-split level dummy, conditional on exposure
\newcommand{\WedgeBreakSpanProfit}{320} % bootstrap 95% CI span on the break date, profit series (days)
\newcommand{\WedgeBreakSpanWedge}{219}  % same, wedge series
\newcommand{\WedgePlaceboSharePct}{38}  % arbitrary split dates reproducing the observed two-window gap
\newcommand{\WedgeZeroSharePre}{36}     % floor slots settling within 0.05 JPY/kWh of the DA price
\newcommand{\WedgeZeroSharePost}{73}
\newcommand{\WedgeMeanPre}{6.65}        % mean floor-slot wedge, JPY/kWh
\newcommand{\WedgeMeanPost}{1.40}
\newcommand{\WedgeMeanDeclinePct}{79}
\newcommand{\WedgePninetyPost}{6.21}        % persistent right tail after the decline
\newcommand{\ChannelPriceK}{+34}      % curtailment-call leg, thousand JPY per MW-yr, FULL sample
\newcommand{\ChannelGenK}{-203}       % generation-forecast leg, full sample (34-203=-169 closes to \CnpeFullK)
\newcommand{\ChannelPricePreK}{+47}   % same legs, earlier window (47-299=-252)
\newcommand{\ChannelGenPreK}{-299}
\newcommand{\ChannelPricePostK}{-2}   % later window (-2+50=+48)
\newcommand{\ChannelGenPostK}{+50}
\newcommand{\ChannelGenAbsK}{203}    % magnitude, for "lowers it by ..." phrasing
\newcommand{\ChannelInteractionK}{+19} % 2x2 interaction: order-dependence of the two legs, full sample
\newcommand{\KAboveCostShare}{0.6}          % share of slots with K above marginal cost
\newcommand{\DayAheadFloorCurtail}{0.020}   % mean day-ahead price, floor + curtailment
\newcommand{\MarginalCostFloorCurtail}{9.18} % mean marginal balancing price there
\newcommand{\MarginalCostFloorCurtailMed}{7.52}
\newcommand{\CostExceedsDaShare}{100}       % ... exceeds the day-ahead price this often
\newcommand{\NineFloorCurtailSlots}{617}    % all nine areas at the floor AND curtailing
\newcommand{\MarginalCostNineFloor}{9.55}   % marginal cost in that tightest state
\newcommand{\MarginalCostBlockLong}{5.65}   % restricted to recorded block surplus
\newcommand{\MarginalCostAreaAlone}{10.32}  % area alone in its block
\newcommand{\ZeroRuleAppliedShare}{100}     % K is exactly zero when the block is long
\newcommand{\BlockLongShareFloorCurtail}{54.5} % ... but the block is long only this often
\newcommand{\BreakSpanConditioned}{254}     % break-date CI span, days, at the
\newcommand{\CnpeFullSampleDays}{547} % scored decision test days, full sample
\newcommand{\CnpeFullK}{-169}       % full-sample CNPE, thousand JPY per MW-yr
\newcommand{\CnpeFullLoK}{-295}
\newcommand{\CnpeFullHiK}{-53}
\newcommand{\CnpeFullAbsK}{169}     % magnitude, for "costs ... thousand JPY" phrasing
\newcommand{\CnpeFullP}{0.001}      % HLN-DM p at L=3, full sample
\newcommand{\FloorPriceRatio}{459}
\newcommand{\DecShareCommitChangePct}{58.4} % share of slots whose committed quantity changes
\newcommand{\DecSumAbsDq}{117.7}        % annualised commitment shift, MWh per MW-installed-yr
\newcommand{\DecSumAbsDqLo}{107.0}
\newcommand{\DecSumAbsDqHi}{128.6}
\newcommand{\DecProfitDiffAbsK}{35.2}   % post-gate commitments settle this much higher, thousand JPY per MW-yr
\newcommand{\DecProfitDiffLoAbsK}{3.8}
\newcommand{\DecProfitDiffHiAbsK}{76.4}
\newcommand{\DecProfitDmPLthree}{0.050}
\newcommand{\DecProfitDmPLnine}{0.057}
\newcommand{\CnpePostRunK}{-134}        % CNPE of the post-gate run vs price-only, full sample
\newcommand{\NetImprovement}{31.3}       % MWh per MW-installed-yr, EXACT marginal change in block |imbalance| (phase77)
\newcommand{\NetImprovementLo}{13.2}
\newcommand{\NetImprovementHi}{48.0}
\newcommand{\GrossReduction}{106.3}      % the same reduction before netting
\newcommand{\SurvivingShare}{29}         % net / gross, per cent, exact
\newcommand{\SurvivingShareFloor}{52}    % on floor slots -- the paper's subject (exact)
\newcommand{\NetImprovementFloor}{11.2}
\newcommand{\NetImprovementAlone}{8.3}   % areas alone in their block (exact)
\newcommand{\NetImprovementDirClassified}{42.9} % the superseded direction-classified figure, kept for the S4 comparison
\newcommand{\NetFlipSlots}{14}          % slots where the marginal portfolio can flip the block's sign
\newcommand{\NetPanelSlots}{204{,}885}
\newcommand{\NetExactLinearGapPct}{0.02} % exact vs first-order-in-B gap, per cent
\newcommand{\NetImprovementPvW}{33.4}   % exact net under PV-capacity weights
\newcommand{\NetImprovementPvWLo}{12.7}
\newcommand{\NetImprovementPvWHi}{52.8}
\newcommand{\NetImprovementPeakW}{44.3} % exact net under peak-demand weights
\newcommand{\NetImprovementPeakWLo}{19.1}
\newcommand{\NetImprovementPeakWHi}{68.5}
\newcommand{\DevSignFlipSharePct}{23.2}  % slots where the op forecast moves the deviation across zero -- why the direction-classified figure overstated
\newcommand{\CrossAreaNettingRatio}{0.772} % share of gross surviving signed summation
\newcommand{\CrossAreaOffsetPct}{23}     % ... so only this much offsets
\newcommand{\CnpeLevelPreK}{+13.3}    % level component, earlier window, k JPY per MW-yr
\newcommand{\CnpeCovPreK}{-265.5}     % covariance component, earlier window
\newcommand{\CnpeLevelPostK}{+28.3}   % level component, later window
\newcommand{\CnpeCovPostK}{+19.6}     % covariance component, later window
\newcommand{\CnpeBandRthreeTwentyK}{-151}  % CNPE under tolerance-band settlement at eps=0.20, thousand JPY per MW-yr
\newcommand{\ToleranceBandReliefPct}{40}   % private penalty removed by tolerance-band settlement at eps=0.20, at zero physical cost
\newcommand{\WedgePostInteraction}{+17.3}   % post x floor-exposure interaction: penalty weaker later
\newcommand{\WedgePostInteractionP}{0.061}
\newcommand{\WedgePostVarSharePct}{2.4}     % share of floor-exposure variance in the later window
\newcommand{\WedgeSlopePre}{-37.5}          % subsample slope, earlier window
\newcommand{\WedgeSlopePost}{-19.0}         % subsample slope, later window (thin base)
\newcommand{\CnpeAbsImbPreAbs}{127}
\newcommand{\CnpeAbsImbPreLoAbs}{155}
\newcommand{\CnpeAbsImbPreHiAbs}{98}
\newcommand{\CnpeAbsImbPostAbs}{60}
\newcommand{\CnpeAbsImbPostLoAbs}{100}
\newcommand{\CnpeAbsImbPostHiAbs}{21}
\newcommand{\CnpePreLoAbsK}{415}
\newcommand{\CnpePreHiAbsK}{109}
\newcommand{\CnpeBandSigEpsPct}{5}      % same as \CnpeBandSigEps, as a percentage
% ===== END paper-B-only ===================================================

% Direct floor models (phase81): the winner of five, with the weather share the
% price-only control isolates. Reported with its multiplicity, per the precommit.
\newcommand{\FloorDirectBest}{0.653}
\newcommand{\FloorDirectIsotonic}{0.597}
\newcommand{\FloorDirectPriceOnly}{0.376}
\newcommand{\FloorDirectWeatherShare}{42\%}

% Floor-event panel (phase66v2): the calibration chain needs a burn-in, so the
% floor panel is shorter than the ladder's and the manuscript says so rather
% than letting the reader assume one number covers both.
\newcommand{\FloorTestDays}{702}
\newcommand{\FloorTestSlots}{302{,}454}

% Second-provider check (phase86, NCEP GFS): external generality, single-point
% arm, same panel and folds as the ladder. Not a headline replacement.
\newcommand{\GfsPreGateImp}{0.203}
\newcommand{\GfsPreGateImpLo}{0.166}
\newcommand{\GfsPreGateImpHi}{0.241}
\newcommand{\GfsJmaMinusGfs}{0.022}
\newcommand{\GfsJmaMinusGfsLo}{0.009}
\newcommand{\GfsJmaMinusGfsHi}{0.036}
\newcommand{\GfsJmaMinusGfsP}{0.0007}
\newcommand{\GfsVintageStep}{0.002}
\newcommand{\GfsVintageStepLo}{-0.003}
\newcommand{\GfsVintageStepHi}{0.008}
\newcommand{\GfsVintageStepP}{0.38}

% ---------------------------------------------------------------------------
% Replication-archive identifiers. Applied Energy's research-data policy asks
% for a deposit at SUBMISSION with the dataset cited and linked in the paper,
% so these are placeholders only until the archive is minted -- which is a
% publication event and waits on PI sign-off. They are EXEMPT from provenance
% binding by construction: no analysis produces them.
% ---------------------------------------------------------------------------
\newcommand{\ReplicationDoi}{10.5281/zenodo.22043820}
% Zenodo's version-independent identifier for the same deposit. Cited
% alongside the version DOI so a reader always reaches the latest
% archived version, while the version DOI and the commit pin exactly
% what produced the reported numbers.
\newcommand{\ReplicationConceptDoi}{10.5281/zenodo.22041072}
\newcommand{\ReplicationCommit}{\texttt{817d351}}

% Audit reliability (phase88b). The raw six-label agreement is moderate; the
% statistic that binds the paper's claims is the Yes-vs-not-Yes flip rate,
% because every reported rate is a share coded Yes.
\newcommand{\LitCoderPapers}{28}
\newcommand{\LitCoderPairs}{624}
\newcommand{\LitAgreement}{75.3}
\newcommand{\LitKappa}{0.65}
\newcommand{\LitYesFlipPct}{1.0}
\newcommand{\LitYesFlipN}{6}
\newcommand{\LitYesFlipKappa}{0.94}

% ---------------------------------------------------------------------------
% Physical target, rebuilt without look-ahead (phase89) and the mechanism test
% (phase90). The MW figures are PER-AREA MEANS over the eight areas whose
% operators publish half-hourly output, not fleet totals.
% ---------------------------------------------------------------------------
\newcommand{\PhysDelta}{59.9}
\newcommand{\PhysDeltaLo}{49.6}
\newcommand{\PhysDeltaHi}{70.6}
\newcommand{\PhysDeltaDays}{696}
\newcommand{\PhysPreGateMae}{878}
\newcommand{\PhysPostGateMae}{819}
\newcommand{\PhysStaleMae}{1322}
\newcommand{\PhysOracleMae}{737}
\newcommand{\PhysLookaheadCost}{8.6}
\newcommand{\PhysStampCost}{134}
\newcommand{\PhysDeltaExcl}{64.6}
\newcommand{\PhysDeltaRecorded}{56.9}
\newcommand{\PhysLatestMae}{828}
\newcommand{\PhysAreas}{eight}
\newcommand{\PhysDeltaWide}{70.1}
\newcommand{\PhysDeltaWideDays}{806}
% Price-only baseline MAE on the sub-panel Table 6 reports: the same number
% that table prints, so the prose and the table cannot disagree.
\newcommand{\PriceBaselineMae}{2.142}
% Absolute-performance panel (phase82b): the sub-panel where the naive 168-hour
% benchmark exists, MEASURED rather than inferred. The filter is on slots, so
% although 14 days lose some benchmark slots, four of them keep the rest and the
% sub-panel spans 722 days -- not the 732 - 14 = 718 an earlier version printed.
\newcommand{\AbsPanelDays}{722}
\newcommand{\AbsPanelSlots}{310{,}284}
% Enriched-baseline decomposition (phase54v2). The interval macros exist because
% two of these were hand-typed from the SUPERSEDED phase54 run and stayed there
% through two revisions; the ratchet cannot see numbers inside prose brackets.
\newcommand{\EnrichedSlots}{265{,}872}
\newcommand{\EnrichGainAdmLo}{0.004}
\newcommand{\EnrichGainAdmHi}{0.036}
\newcommand{\ProxyCarriedLo}{0.159}
\newcommand{\ProxyCarriedHi}{0.223}
\newcommand{\EnrichPerAreaLo}{0.087}
\newcommand{\EnrichPerAreaHi}{0.289}
\newcommand{\AbsNaiveMae}{3.096}
\newcommand{\AbsPreGateMae}{1.890}
\newcommand{\AbsPreGateRmae}{0.611}
\newcommand{\AbsFloorNaiveMae}{5.481}
\newcommand{\AbsFloorBaselineMae}{5.223}
\newcommand{\AbsFloorPreGateMae}{3.383}
\newcommand{\FloorAbsGain}{1.84}
% Floor-model comparison under paired inference (phase81b stage 1). The margins
% are differences of the MACRO statistic the text reports, from a block
% bootstrap over delivery days with all contrasts sharing one set of draws.
\newcommand{\FloorWinnerMargin}{0.056}
\newcommand{\FloorWinnerMarginLo}{0.030}
\newcommand{\FloorWinnerMarginHi}{0.080}
% The bootstrap p is AT its resolution floor (1/10,000) for all five contrasts,
% so the Holm-adjusted value is an upper bound set by the draw count, not a
% measured p. S4 says so.
\newcommand{\FloorWinnerHolmP}{0.0008}
\newcommand{\FloorWeatherMargin}{0.276}
\newcommand{\FloorWeatherMarginLo}{0.217}
\newcommand{\FloorWeatherMarginHi}{0.327}
\newcommand{\FloorGateContrast}{0.015}
\newcommand{\FloorGateContrastLo}{-0.001}
\newcommand{\FloorGateContrastHi}{0.030}
% The floor-model comparison AFTER the tie artefact was found. Scoring the
% isotonic chain by its calibrated probability charges it for the ties its own
% monotone calibrator creates (19,898 distinct values against 284,101 for the
% ranking it induces), and two thirds of the apparent PR-AUC margin was that.
% On the like-for-like ranking the two model classes are indistinguishable; the
% log score is the instrument that separates them, which is the point.
\newcommand{\FloorRankMargin}{0.018}
\newcommand{\FloorRankMarginLo}{-0.004}
\newcommand{\FloorRankMarginHi}{0.039}
\newcommand{\FloorRankMarginP}{0.114}
\newcommand{\FloorIsotonicRank}{0.635}
\newcommand{\FloorBrierIsotonic}{0.0184}
\newcommand{\FloorBrierDirect}{0.0180}
\newcommand{\FloorBrierMarginP}{0.208}
\newcommand{\FloorLogIsotonic}{0.0725}
\newcommand{\FloorLogDirect}{0.0648}
\newcommand{\FloorLogMargin}{-0.0077}
\newcommand{\FloorLogMarginLo}{-0.0116}
\newcommand{\FloorLogMarginHi}{-0.0040}
\newcommand{\FloorGateContrastDirect}{-0.001}
\newcommand{\FloorGateContrastDirectLo}{-0.011}
\newcommand{\FloorGateContrastDirectHi}{0.013}
\newcommand{\FloorIsotonicTiesCal}{19{,}898}
\newcommand{\FloorIsotonicTiesRaw}{284{,}101}
\newcommand{\FloorStageTwoGap}{0.011}
\newcommand{\FloorDirectBestMicro}{0.692}
\newcommand{\FloorDirectIsotonicMicro}{0.656}
\newcommand{\FloorDirectBestSkill}{0.641}
\newcommand{\FloorDirectIsotonicSkill}{0.584}
% Floor positives per area: the caption said "a few hundred" where the range is
% 350 to 2,430 -- understating most areas by up to sevenfold. Bound by aggregate
% so the area names in the sentence are asserted too.
\newcommand{\FloorSlotsAreaLo}{350}
\newcommand{\FloorSlotsAreaHi}{2{,}430}
% The price target ON THE COMMON PANEL (phase89b). Figure 5's panel (a) used to
% show the full-panel ladder while its panel (b) showed the common panel, and
% the caption claimed both were the common panel.
\newcommand{\CommonPriceImpPre}{0.259}
\newcommand{\CommonPriceImpPost}{0.261}
\newcommand{\CommonPriceImpStale}{0.045}
% Physical-target construction (phase89 .conventions / .metadata).
\newcommand{\CapVintage}{2024-09}
% Inclusive hour filter [9, 15] over half-hourly slots labelled by start:
% fourteen slots per area-day, 09:00 through 15:30. It read 09:00--15:00
% until 2026-08-13, understating the window by an hour.
\newcommand{\PhysWindow}{09:00--15:30}
\newcommand{\PhysWindowSlots}{fourteen}
\newcommand{\PhysMinTrain}{200}
\newcommand{\PhysBootDraws}{10{,}000}
\newcommand{\PhysBootBlock}{seven}
\newcommand{\PhysPerAreaLo}{164}
\newcommand{\PhysPerAreaHi}{2{,}353}
% Mechanism test: shares of residual VARIANCE explained by the innovation.
% The two placebo regressions are a pre-registered VOID condition: had either
% been significant, every slope in the mechanism table would have had to be read
% as an artefact of the estimator. Their p-values were not bound to anything --
% the one quantity that can void the table was the one nothing checked. The
% physical placebo sits at 0.051, just inside the threshold, which the supplement
% now states rather than leaving a reader to notice.
\newcommand{\MechPlaceboPhysP}{0.051}
\newcommand{\MechPlaceboPriceP}{0.321}
\newcommand{\MechPhysR}{14.0}
\newcommand{\MechPriceR}{1.2}
\newcommand{\MechPriceSlope}{-0.39}
\newcommand{\MechInnovSd}{531}
\newcommand{\MechSlots}{77{,}744}
\newcommand{\MechClusters}{5{,}560}
% What the innovation could be worth in price units. TWO panels, and the
% distinction matters: 0.011 is the ceiling on the slots where the innovation
% exists (daylight, eight areas, pre-gate residual MAE 2.45), while 0.003 is the
% same thing expressed on the panel the ladder scores, which is what the
% equivalence margin and the gate-timing delta are measured on. An earlier
% version compared the sub-panel figure with full-panel quantities and called the
% ceiling "about one tick"; like-for-like it is about a third of a tick.
\newcommand{\MechImpliedMae}{0.011}
\newcommand{\MechImpliedTicks}{1.13}
\newcommand{\MechImpliedMaeFull}{0.004}
% Four-decimal forms. Section 5 introduces these as "three numbers in a revealing
% order" and then printed the first two as 0.003 and 0.003: at three decimals the
% order the sentence promises is invisible, and the stated 75% reads as 100%.
% The same collision was fixed in the Conclusions and in Section 7 by quoting the
% share instead; here the ordering IS the argument, so the precision goes up.
\newcommand{\MechCeilingPrecise}{0.0036}
\newcommand{\MechCeilingOls}{0.0034}   % what a least-squares projection reaches
\newcommand{\MechImpliedTicksFull}{0.36}
\newcommand{\MechInnovSlotShare}{24.6}
% Capacity-weight sensitivity (phase91). The three arms differ only in the point
% weights; the 2024-09 arm reproduces the published ladder exactly, which is what
% licenses reading the other two as weight effects.
\newcommand{\WeightHeadlineBase}{0.2529}
\newcommand{\WeightHeadlineLater}{0.2508}
\newcommand{\WeightHeadlineEqual}{0.2539}
\newcommand{\WeightLaterVintage}{2025-12}
\newcommand{\WeightHeadlineSpread}{0.0031}
% The pre-registered one-tick equivalence test under each weighting. Gate 59
% landed on branch b BECAUSE of these: the claim holds on the paper's weights and
% with the register discarded, and fails under the later snapshot. An earlier
% version of the gate tested only the point estimate's sign and CI overlap, and
% would have passed while the claim it exists to protect failed.
\newcommand{\WeightEquivPBase}{0.023}
\newcommand{\WeightEquivPLater}{0.077}
\newcommand{\WeightEquivPEqual}{0.005}
\newcommand{\WeightGateDeltaLater}{0.0049}
